% !TEX root = ../../../tree_notes.tex
% =====================================================================
%  12.4  Eccentricity, Radius, Centre (+ reroot preview)  sec:c09-radius
%  Exports: def:c09-ecc  thm:c09-eccend  thm:c09-centerrule
%           lst:c09-eccsweep  fig:c09-eccprof  ex:c09-ecc-run
%           prob:c09-capital  lst:c09-219d-soln   (practice: LC 317)
%  Solved slot: CF 219D Choosing Capital for Treeland (1700) -- the
%  live reroot-shift preview whose weighted sibling is the Summit.
%  (The plan's first pick, CF 592D Super M, verified at 2200 -- demoted
%  to the mixed set to keep the solved ladder below Summit 1900.)
%  LC 310 minimum-height-tree peeling stays Ch.4's -- referenced, not
%  re-taught. Distances here are one-shot; per-query stays Ch.16.
% =====================================================================
\section{Eccentricity, Radius, Centre --- and the Best Root}\label{sec:c09-radius}

\begin{previouslybox}
\begin{enumerate}[nosep, leftmargin=1.6em]
  \item \thmref{c09-twosweep} ends with two distance arrays in hand (the sweeps'
        $\mathit{dist}[]$). What does the \emph{second} array's maximum compute? ($D$
        itself --- and its per-node values are distances from a diameter \emph{end})
  \item On the diametral path of $T$, node~$1$ sits $3$ steps from end~$8$ and $4$ from
        end~$10$. Which is $\mathit{ecc}(1)$ --- and why is checking just those two
        nodes enough? (max $=4$; this section's \thmref{c09-eccend} is the reason)
  \item 12.3 warned the centroid $\ne$ the centre and named the broom's two answers.
        Which one is defined by \emph{counts}, which by \emph{distances}? (centroid by
        counts, centre by distances)
  \item \coderef{c09-dfs3} computed $\sum_u d(1,u) = 20$ for $T$ via $\depth$. Suppose
        you knew the answer for root~$1$; what single edge-flip information would give
        the answer for a child root? (how many nodes are on each side of the edge ---
        $n - \size(v)$ get farther, $\size(v)$ get nearer)
\end{enumerate}
\end{previouslybox}

\begin{topicintro}
Three questions share one geometry. \emph{How far is $v$ from the worst node?} ---
the \term{eccentricity} $\mathit{ecc}(v)$. \emph{What is the best possible worst-case,
and where?} --- the \term{radius} $R = \min_v \mathit{ecc}(v)$ and the \term{centre},
the argmin set. These are the ``put the warehouse at the node minimising the maximum
delivery distance'' answers --- and unlike almost every other distance problem in
this book, they cost \emph{two sweeps total}, not $n$ sweeps, because of a single
structural identity: \emph{the farthest node from anywhere is always one of the two
diameter ends} (\thmref{c09-eccend}). With the diameter ends $A, B$ from
\secref{c09-diameter} and their two distance arrays, the whole table
$\mathit{ecc}[\cdot]$ is a $\max$ per node, the centre falls out as the argmin --- and
\thmref{c09-centerrule} upgrades ``found it'' to ``it is literally the middle of every
diameter,'' the exact fact behind the leaf-peeling solution of LC~310's
minimum-height-trees family (rooted treatment in Chapter~4, \chref{c04}; we keep the
distance side).

The section then turns from $\max$ to $\sum$: replace ``farthest'' by ``sum of all
distances'' and the identity dies --- but a weaker, mechanical one lives: the
\term{reroot shift}, $f(v) = f(u) + S - 2\,\size(v)$, one DFS plus one replay to get
every root's total at once. That preview is the entire budget this chapter spends on
Chapter~15 (\chref{c20}); what it buys immediately is the solved problem, Codeforces~219D
\emph{Choosing Capital for Treeland} (\probref{c09-capital}) --- one-way roads, choose
the capital minimising forced reversals --- and its weighted, full-budget sibling is
the chapter-end Summit, CF~1092F.
\textbf{Prerequisites.} \secref{c09-diameter} (ends $A,B$, sweeps, the branch-depth
claim from its hardening box); \secref{c09-measures} ($\size$, and \coderef{c09-iter}
discipline). Per-pair \emph{query} distance remains Chapter~16's (\chref{c11}).

\textbf{After this section you can}
\begin{itemize}[nosep]
  \item compute all eccentricities, the radius, and the centre from two sweeps; \hfill\lv{3}
  \item prove $\mathit{ecc}(v)=\max(d(v,A),d(v,B))$ and the centre-is-the-middle rule; \hfill\lv{4}
  \item run the unweighted and reversal-cost reroot shifts over all roots in $\Theta(n)$; \hfill\lv{4}
\end{itemize}
\end{topicintro}

\subsection{Intuition}

Picture the diametral path $A \rightsquigarrow B$ as a road, and every other node as a
house on a spur off some milepost. The branch claim proved by example in
\secref{c09-diameter}'s hardening box says a spur can never run longer than the nearer
road-end it hangs off --- otherwise the road itself wasn't the diameter. Now stand
anywhere, at node~$v$. Look at any target $u$: getting from $v$ to $u$ means climbing
out to the road, sliding along it, and maybe taking $u$'s spur. But $u$'s spur is short
(the claim), so sliding all the way to \emph{whichever road-end is farther from your
milepost} always reaches at least as far. The worst node from $v$ is, up to spur
detours, one of the two ends --- and taking $\max(d(v,A), d(v,B))$ never underestimates,
because $A$ and $B$ \emph{are} nodes. One identity, two sweeps, whole table.

The centre then reads itself off: walk the diameter road and your worst-case distance
at milepost $i$ is $\max(i, D-i)$ --- two arms stretching in opposite directions, each
growing by one as you walk. The minimum of that V-shaped function is where the arms
cross, $i = D/2$: one node when $D$ is even, the two middle nodes of an odd $D$ ---
which is why $T$ ($D=7$) has \emph{two} centres. Contrast 12.3 one final time: the
centroid hiked by \emph{headcounts}, the centre by \emph{mileposts}; on $T$ both
happen to answer $\{1,3\}$ --- a coincidence the broom breaks instantly.

The $\sum$-version needs a different machine: no farthest-end shortcut exists for
averages, only a bookkeeping law. Cross edge $u \to v$ with root: the $\size(v)$ nodes
inside $v$'s side each step one unit \emph{closer}, the other $n - \size(v)$ each step
one \emph{farther} (weights: replace counts by mass $W_v$ vs $S - W_v$) --- so every
root's answer is the parent's answer plus a two-term correction. That is one
$\Theta(n)$ replay for \emph{all} roots, and it is the whole of \probref{c09-capital}
once ``distance'' becomes ``number of wrongly-one-way roads.''

\subsection{Definitions and the Two Structure Theorems}

\begin{definition}{Eccentricity, radius, centre}{c09-ecc}
$\mathit{ecc}(v) = \max_{u \in V} d(v, u)$; the \term{radius} is
$R = \min_v \mathit{ecc}(v)$; the \term{centre} is the set of nodes attaining $R$.
(Compare: the \term{periphery} is the set attaining $\max_v \mathit{ecc}(v) = D$ ---
the diameter ends are peripheral, and \lemref{c09-sweepend} said every sweep finds one.)
A single-node tree: $\mathit{ecc} = R = 0$ and the centre is that node.
\end{definition}

\begin{theorem}{Two endpoints suffice for every eccentricity}{c09-eccend}
Let $A \rightsquigarrow B$ be any diametral path. Then for every node $v$,
\[
  \mathit{ecc}(v) \;=\; \max\big(d(v, A),\; d(v, B)\big),
\]
computable for all $v$ from the two sweeps of \coderef{c09-sweep} in $\Theta(n)$.
\end{theorem}
\begin{proofbox}
$(\ge)$ is instant: $A,B \in V$. For $(\le)$ let $u$ be any target; write $p$ for the
node where $v$'s road meets $P$ (the unique node of $P$ on the path from $v$ to~$P$;
$v$ itself counts, $p = v$, when $v$ is on the road), and $q$, $t_u$ analogously for
$u$. The branch claim (\secref{c09-diameter}, hardening item~2: a spur of depth $t$ at
road-position $q$ would extend the diameter unless $t \le \min(d(q,A), d(q,B))$) gives
$t_u \le \min(d(q,A), d(q,B))$, and then
\[
  d(v,u) = \underbrace{d(v,p)}_{\text{out}} + d(p,q) + t_u
  \;\le\; d(v,p) + \min\big(d(p,q){+}d(q,A),\; d(p,q){+}d(q,B)\big).
\]
Whichever side of $p$ the point $q$ lies on, one sum collapses exactly: $q$ toward $B$
makes $d(p,q){+}d(q,B) = d(p,B)$; $q$ toward $A$ makes $d(p,q){+}d(q,A) = d(p,A)$. So
the $\min$ is at most $\max(d(p,A), d(p,B))$ no matter the order, and
$d(v,u) \le d(v,p) + \max(d(p,A), d(p,B)) = \max(d(v,A), d(v,B))$. Maximising over $u$
proves $(\le)$.
\end{proofbox}

\begin{theorem}{The centre is the (half-)middle of every diameter}{c09-centerrule}
The centre of a tree consists of the unique node at distance $D/2$ from both ends when
$D$ is even, or the two adjacent nodes flanking the middle \emph{edge} when $D$ is
odd; in particular $R = \lceil D/2 \rceil$, the centre is the same for \emph{all}
diametral paths, and it is one node or an adjacent pair --- never more.
\end{theorem}
\begin{proofbox}
By \thmref{c09-eccend}, along the path $A \rightsquigarrow B$ the node at position $i$
has $\mathit{ecc} \ge \max(i, D-i)$, with equality \emph{at} those positions because
$A, B$ are witnesses. The function $g(i) = \max(i, D-i)$ is strictly decreasing up to
$i = D/2$ and strictly increasing after it: if $D$ is even its unique minimum is
$g(D/2) = D/2$; if odd its minimum $g\big(\tfrac{D \pm 1}{2}\big) = \tfrac{D+1}{2}$ is
attained exactly at the two central nodes, which are adjacent. So the centre of $T$
restricted to \emph{this} path is that one node or pair, at radius $\lceil D/2 \rceil$.
No node off the path can join the centre: at road-position $p$ with spur depth $t \ge 1$,
\[
  \mathit{ecc}(\text{spur end}) = \max\big(t + d(p,A),\; t + d(p,B)\big)
  = t + \max(d(p,A), d(p,B))
\]
\noindent
\hspace*{2em}$>\; g(\text{position of } p) \;\ge\; R$ --- factoring $t$ out of the
$\max$; the $g$ step is \thmref{c09-eccend} again (distances from $p$ to the ends are
its two arms along the road). Finally ``same for every diameter'': two
centres would sit at the middle of \emph{both} paths --- and the centre set is a
connected single node or edge by the strictness of $g$, so all middles coincide.
\end{proofbox}

\begin{lstlisting}[caption={The full metric readout: two sweeps -> ecc table -> radius, centre, and the distance-sum reroot}, label={lst:c09-eccsweep}]
// sweep() of lst:c09-sweep, extended to also hand back its dist[] array
// (add a 4th field  vector<int> darr;  or keep the array global):
Sweep S1 = sweep(0);              // any start -> reach an extremity S1.far = a
Sweep S2 = sweep(S1.far);         // from end a: S2.darr[] = d(a, .), far = end b
int D = S2.dist;                  // the diameter
Sweep S3 = sweep(S2.far);         // from end b: S3.darr[] = d(b, .)
int R = n;  vector<int> centre;
for (int u = 0; u < n; u++) {
    int e = max(S2.darr[u], S3.darr[u]);     // ecc(u): max over BOTH ends
    if (e < R) { R = e; centre = {u}; }
    else if (e == R) centre.push_back(u);    // centre: 1 node or an adjacent pair
}
// R == (D+1)/2 (integer ceil); centre == middle of the S2.par walk-back path.
// ---- reroot preview (sums, not maxes): ----
// f[0] = sum_u dep[u]  (one dfs3); then for tree edge u->v (u parent):
//     f[v] = f[u] + n - 2*sz[v];             // one forward replay, all roots
// weighted: f[v] = f[u] + S - 2*W[v] with W = subtree weight sum.
\end{lstlisting}

\begin{complexitybox}
Three sweeps + a scan: $\Theta(n)$, and only \emph{two} sweeps are load-bearing (the
third from $B.\mathit{far}$ is just $\mathit{dist}(A,\cdot)$ --- or reuse the first
sweep's array if it started at $A$... it didn't; three is the honest count, still
linear). The reroot block is one \coderef{c09-dfs3} plus one forward replay,
$\Theta(n)$, arrays only --- no recursion risk with the \coderef{c09-iter} shape.
Verdict: eccentricity/radius/centre and all-root distance sums for $\Theta(n)$ total;
compare the naive $n$ BFSs, $\Theta(n^2)$ --- the identity \thmref{c09-eccend} is what
pays the entire bill.
\end{complexitybox}

\subsection{Worked Example}

\begin{example}{The six-angle table closes: $\mathit{ecc}$, radius, centre of $T$}{c09-ecc-run}
Diameter ends $A = 8$, $B = 10$ (\figref{c09-diampath}). Two distance arrays give every
row of the table below --- each $\mathit{ecc}$ is just the $\max$ of its column pair
(\thmref{c09-eccend}); verify two rows by brute force to feel the identity working:
\[
\begin{array}{c|cccccccccc}
u                & 1 & 2 & 3 & 4 & 5 & 6 & 7 & 8 & 9 & 10 \\ \hline
d(u, 8)          & 3 & 2 & 4 & 3 & 1 & 5 & 5 & 0 & 6 & 7 \\
d(u, 10)         & 4 & 5 & 3 & 6 & 6 & 4 & 3 & 7 & 1 & 0 \\
\mathit{ecc}(u) & 4 & 5 & 4 & 6 & 6 & 5 & 5 & 7 & 6 & 7
\end{array}
\]
$R = \min \mathit{ecc} = 4$, centre $= \{1, 3\}$ --- exactly the two middle nodes of
$8\!-\!5\!-\!2\!\underline{\mid}\!1\!-\!3\!\underline{\mid}\!7\!-\!9\!-\!10$
($D = 7$ odd, adjacent pair, $R = \lceil 7/2 \rceil = 4$ \thmref{c09-centerrule}).
Periphery ($\mathit{ecc} = D = 7$): $\{8, 10\}$ --- and node~$5$? $\mathit{ecc}(5)=6$:
nearly peripheral, never a sweep's prize.\\
\emph{Distance sums (reroot preview).} $f(1) = \sum \depth = 20$. Crossing $1 \to 2$
($\size(2)=4$): $f(2) = 20 + 10 - 2{\cdot}4 = 22$; then $2 \to 5$ ($\size 2$):
$f(5) = 22 + 10 - 4 = 28$; $1 \to 3$ ($\size 5$): $f(3) = 20 + 10 - 10 = 20$;
$3 \to 7$ ($\size 3$): $f(7) = 20 + 10 - 6 = 24$; $7 \to 9$ ($\size 2$):
$f(9) = 24 + 10 - 4 = 30$; $9 \to 10$ ($\size 1$): $f(10) = 30 + 10 - 2 = 38$.
Brute-force check on $f(10)$: distances from $10$ are
$[4,5,3,6,6,4,3,7,1,0]$ $= 38$ \checkmark. Minimum sum at node~$1$ or~$3$ (tie!) ---
$\sum$-optimisers are the \emph{median} family, and here they sit right on the
centroids: not an accident, but a theorem for Chapter~15 (\chref{c20}). \end{example}

\begin{figure}[H]
\centering
\begin{tikzpicture}[diagram, scale=0.9, transform shape,
    tnode/.style={circle, draw=primaryblue, line width=0.9pt, fill=primaryblue!8,
                  minimum size=5.6mm, font=\small\bfseries, text=primaryblue!60!black},
    cen/.style={tnode, draw=softpurple, fill=softpurple!22, line width=1.2pt,
                text=softpurple!60!black}]
  \node[tnode] (v8) at (0,0) {$8$};
  \node[tnode] (v5) at (1.6,0) {$5$};
  \node[tnode] (v2) at (3.2,0) {$2$};
  \node[cen]   (v1) at (4.8,0) {$1$};
  \node[cen]   (v3) at (6.4,0) {$3$};
  \node[tnode] (v7) at (8.0,0) {$7$};
  \node[tnode] (v9) at (9.6,0) {$9$};
  \node[tnode] (v10) at (11.2,0) {$10$};
  \draw[tedgehl] (v8)--(v5)--(v2)--(v1)--(v3)--(v7)--(v9)--(v10);
  \foreach \nd/\i/\e in {v8/0/7, v5/1/6, v2/2/5, v1/3/4, v3/4/4, v7/5/5, v9/6/6, v10/7/7} {
    \node[font=\tiny\ttfamily, text=darkgray, anchor=north] at (\nd.south) {\i};
    \node[font=\tiny\ttfamily, text=warningred!75!black, anchor=south] at (\nd.north) {\e};}
  \node[annot, font=\scriptsize, anchor=north west] at (0,-0.7)
    {position $i$ along the diameter (bottom, grey)};
  \node[annot, font=\scriptsize, anchor=north west, text=warningred!70!black] at (6.6,-0.7)
    {$\mathit{ecc} = \max(i, 7{-}i)$: the V, minimal at the two middles};
\end{tikzpicture}
\caption{The diameter of $T$ stretched straight, carrying its $\mathit{ecc}$ profile
(red, above) and position index (grey, below): $7,6,5,\mathbf{4},\mathbf{4},5,6,7$ ---
$\max(i, D-i)$ in the shape \thmref{c09-centerrule} predicts. The purple pair $\{1,3\}$
is the centre; the V never dips below $R = 4$, and off-road nodes (undrawn) can only
lift it by their spur depth.}
\label{fig:c09-eccprof}
\end{figure}

\begin{tracebox}[Reading \thmref{c09-eccend} against brute force at $v = 4$]
Formula: $\max(d(4,8), d(4,10)) = \max(3, 6) = 6$. Brute: distances from $4$ to
$\{1{:}2, 2{:}1, 3{:}3, 5{:}2, 6{:}4, 7{:}4, 8{:}3, 9{:}5, 10{:}6\}$ --- max $= 6$
at node~$10$, a diameter end \emph{exactly} as the theorem promises. The spur depth of
$4$ is $1$ at milepost $2$: $1 \le \min(2, 5)$ --- the branch claim holding with room
to spare, and $\mathit{ecc}(4) = 1 + \max(2, 5)$: spur out, then all the way to the
far end. \checkmark
\end{tracebox}

\subsection{The Gallery: One Tree, Eight Numbers}

\begin{center}
\small
\renewcommand{\arraystretch}{1.25}
\begin{tabular}{@{}p{0.21\linewidth} p{0.20\linewidth} p{0.24\linewidth} p{0.26\linewidth}@{}}
\toprule
\textbf{Number} & \textbf{Definition} & \textbf{On $T$} & \textbf{Machinery} \\
\midrule
$\depth(v)$ & root-to-$v$ edges & $\depth(10)=4$ (root $1$) & down-current of 12.1 \\
$h(v)$ & deepest node below & $h(1)=4$ & up-current $\max+1$ \\
$\size(v)$ & subtree headcount & $\size(3)=5$ & up-current sum \\
$D$ & longest path & $7$, ends $8,10$ & two sweeps / arm join \\
$\mathit{ecc}(v)$ & farthest from $v$ & $\mathit{ecc}(2)=5$ & $\max(d(v,A),d(v,B))$ \\
$R$ & best eccentricity & $4$ & argmin of the row above \\
centre & argmin set & $\{1,3\}$ & middle of every diameter \\
centroid & smallest max-piece & $\{1,3\}$ (tie of a different kind) & $\size$ + upward piece \\
\bottomrule
\end{tabular}
\end{center}

\subsection{Variations and Twists}

\begin{itemize}[nosep]
  \item \textbf{Absolute vs central vs peripheral 1-median:} $\max$ over $u$ is the
        centre, $\sum$ over $u$ is the \emph{median} (sits at a centroid for trees ---
        check on $T$: $f$-minimum at $1$ or $3$, both centroids), and the farthest pair
        is the diameter. Same distance matrix, three different optimisations; the
        reroot-shift machinery answers the $\sum$ one for all roots at once. \hfill\lv{4}
  \item \textbf{Eccentricity of every node in $\bigO{n^2}$ $\to$ $\bigO{n}$:} the naive
        ``BFS from each node'' becomes three sweeps via \thmref{c09-eccend}; the same
        identity upgrades weighted trees \emph{only} if every sweep's distance pass is
        weighted (positive weights: DFS accumulation suffices on trees). \hfill\lv{4}
  \item \textbf{Adding a leaf, dynamically:} centre and centroid both move along the
        diameter/descent paths at speed $\le 1$ per insertion --- the intuition behind
        CF~342E-style maintenance \emph{after} Chapter~24 arms you. \hfill\lv{5}
  \item \textbf{Two-sweep radius shortcut:} $R = \lceil D/2 \rceil$ needs only $D$ ---
        but the \emph{location} of the centre needs the path, hence \code{par[]} from
        sweep two (\secref{c09-diameter}'s path recovery), walk to position
        $\lfloor D/2 \rfloor$. \hfill\lv{3}
\end{itemize}

\subsection{Common Pitfalls}

\begin{warningbox}[Trap 1 -- computing ecc(v) as height(v)]
Under root~$1$, $h(1) = 4 = \mathit{ecc}(1)$ --- pure coincidence of the rooting.
$\mathit{ecc}(v)$ measures \emph{all} directions; $\height(v)$ only down. On $T$ rooted
at~$1$: $\mathit{ecc}(2) = 5$ but $h(2) = 2$ (the farthest node from $2$ is $10$, above
and away). The \emph{tree} radius $\ne$ ``$\height$ of a poor rooting''; get the rooting
out of your head and the two ends into your arrays (\thmref{c09-eccend}).
\end{warningbox}

\begin{warningbox}[Trap 2 -- centre/centroid conflated]
The insight box broom: centroid $=2$, centre $=3$. On $T$ both are $\{1,3\}$ --- and a
solution that passes one by passing the other is a solution on borrowed luck.
``Balance'' problems (halves, recursing) want \emph{sizes}; ``worst-case distance''
problems (facilities, height) want \emph{mileposts}. The first question of design is
which of the two the statement optimises.
\end{warningbox}

\begin{warningbox}[Trap 3 -- one end instead of both]
$\mathit{ecc}(v) \ne d(v, A)$ alone --- $A$ is far from one \emph{side} of the tree.
Node~$8$ with only-end-$A$ reasoning reads $d(8,8) = 0$: ``perfectly central,'' while
it is peripheral at $7$. The identity is a $\max$ over \emph{both} ends; a
single-sweep-from-$A$ table answers ``distance to $A$'', nothing else. (Symptom: your
radius comes out $0$ or near-$0$ on any tree with a long diameter, because an endpoint
is always in the array you forgot.)
\end{warningbox}

\subsection{Problem Lab}

The anchor is Codeforces~219D, \probref{c09-capital}: the reroot shift as a
\emph{cost} version --- no distances at all, just counts of wrongly-pointed edges, and
the same two-phase $\Theta(n)$ machine as \coderef{c09-eccsweep}'s last block. Its
weighted-distance sibling is the Summit.

% ---------------------------------------------------------------------------
\solvedproblem{c09-capital}{Choosing Capital for Treeland}{Codeforces 219D}
  {https://codeforces.com/problemset/problem/219/D}{CF 1700 · L4}
  {reroot shift, directed-edge costs, two-phase DFS, output all argmins}

\begin{problemstatement}[Choosing Capital for Treeland (statement, samples and rating verified against the live problem)]
Treeland has $n$ cities connected by $n-1$ \emph{one-way} roads which, if every road
were bidirectional, would form a tree. Choose a capital: from it, all other cities must
be reachable, reversing as few roads as necessary. Print the minimum number of
reversals on the first line, and \emph{all} capital cities achieving it in ascending
order on the second. Constraints: $2 \le n \le 2\cdot 10^5$; roads given as pairs
$(a, b)$ meaning direction $a \to b$.
\textbf{Sample~1}: $n=3$, roads $2\!\to\!1$, $2\!\to\!3$ $\Rightarrow$ \texttt{0} then
\texttt{2} (city~$2$ already points at both). \textbf{Sample~2}: $n=4$, roads
$1\!\to\!4$, $2\!\to\!4$, $3\!\to\!4$ $\Rightarrow$ \texttt{2} then \texttt{1 2 3} (a
leaf-capital keeps its own outward road and flips the other two; capital~$4$ would flip
all three).
\end{problemstatement}

\begin{attemptgate}[Attempt it yourself first -- \lv{4}, give yourself 25 minutes]
The ``distance'' here is a count of rebellious arrows --- and moving the capital across
one road changes exactly \emph{one} road's compliance.
\peek{Root the (undirected) tree at city~1. What does a road oriented $b \to a$ cost
when the root's flow wants $a \to b$?}
\nudge{One reversal. So $f(1) = \#$roads pointing toward the flow's opposite. Now move
the root across a single edge $\{u,v\}$: every other road keeps its compliance
unchanged --- so $f(v) - f(u)$ depends on \emph{one} edge only.}
\fullhint{Give traversal edges costs: walking $a \to b$ on a given road $a \to b$ costs
$0$, walking against costs~$1$. One DFS from~$1$ sums $f(1)$; then the shift along tree
edge $u \to v$: if the road was $u \to v$, moving the root \emph{costs} $+1$ (that edge
now reads wrong); if $v \to u$, it \emph{saves} $1$: $f[v] = f[u] + 1$ or $f[u] - 1$
respectively --- $f[v] = f[u] + 1 - 2w$ with $w \in \{0,1\}$ the down-cost. Forward
replay, take the min, list argmins ascending (nodes visited in index order gives that
free). $n$ up to $2\cdot 10^5$: iterative two-phase, \coderef{c09-iter} shape.}
\end{attemptgate}

\solstep{Understanding the Problem}
Every root $r$ orients all $n-1$ edges away from $r$ (reachability = the rooted arborescence);
an edge resists iff its given direction points \emph{toward}~$r$'s side. The score
$f(r)$ is a per-edge \emph{sum}, not a $\max$ --- \thmref{c09-eccend} is dead here, the
reroot shift is alive. Output demands \emph{all} argmins sorted, so the second phase
must visit nodes in label order (or sort at the end; ascending scan is free).

\solstep{Key Observations}
\begin{observation}
Rooting at $v$ instead of its tree-neighbour $u$ flips the orientation requirement of
\emph{exactly one} edge --- $\{u,v\}$ itself. All other edges keep the same side-relationship
to the root. Hence consecutive roots differ by $\pm 1$:
$f[v] = f[u] + (1 - 2w)$, $w = $ cost of the required $u \to v$ direction ($0$ if the
road exists that way, $1$ if against). \end{observation}
\begin{observation}
This is the distance-sum shift $f[v] = f[u] + n - 2\size[v]$'s \emph{binary-cost}
sibling: replace counts by per-edge costs and the $\pm$ becomes $\pm 1$ weighted by the
single flipped edge --- the same one-DFS-plus-replay budget, $\Theta(n)$.
\end{observation}
\begin{observation}
With all roads pointing \emph{into}~$4$ (Sample~2), $f(4) = 3$ and each leaf capital
saves its own edge but inherits the other direction wrong: $f(\text{leaf}) = 3 - 1 = 2$
-- the shift runs \emph{downward} from any computed root, direction of bookkeeping
arbitrary.
\end{observation}

\solstep{Think It Through With Me}
\begin{thinkbox}
\deadend{Try every capital: for each, BFS/DFS the directed graph counting reversals.}
{$\bigO{n^2}$ ($2\cdot 10^5$ roots $\times$ linear scans $=4\cdot 10^{10}$) --- TLE;
and every root's answer is the previous one's answer $\pm$ one edge's vote, so $99.999\%$
of the recomputation is already known.}
\deadend{Count in-degrees and pick the ``most outward-pointing'' node.}
{Local degree does not see global orientation: a long chain $1\to2\to3\to\cdots$ has
$f(1) = 0$ but every non-endpoint also has one in, one out --- the answer needs a
\emph{tree-wide} quantity propagated, not a per-node count.}
\aha{Compute $f$ for one root in $\Theta(n)$, then the single-edge shift
$f[v] = f[u] + 1 - 2w$ along the tree: one forward replay produces every root's cost.
The preview paragraph of this section, with arrows for roads.}
\end{thinkbox}

\solstep{Solution}
Store, for each given road $a \to b$, both traversals: $a \to b$ with cost~$0$,
$b \to a$ with cost~$1$ (the cost of \emph{disobeying} the sign). Phase~1: preorder
from node~$0$, accumulating $f[\text{root}] = \sum$ costs of the chosen parent$\to$child
traversals. Phase~2: walk $\mathit{order}$ forward; for each node $v \ne$ root with
incoming down-cost $w_v$ (cost of $par[v] \to v$): $f[v] = f[par[v]] + 1 - 2 w_v$.
Min over $f$; print it, then all $v$ with $f[v] = \min$ in label order.

\solstep{Proof of Correctness}
\begin{proofbox}[The shift is exact, hence argmins are all capitals]
Fix root $u$ and a tree edge $\{u,v\}$ oriented $u \to v$ in the $u$-rooting. Rooting at
$v$ instead: the required orientation flips \emph{only} for $\{u,v\}$ --- for any other
edge $\{x,y\}$, the side of the split it sits on relative to the root is unchanged
(the unique path from the root to any node $w$ changes by prepending $v \to u$ but the
\emph{directions} of all other edges on all paths stay the same), so its compliance is
unchanged. $\{u,v\}$ itself: compliant under root $u$ iff the road reads $u \to v$
(cost $w=0$); under root $v$ it is required as $v \to u$, compliant iff $w = 1$.
Hence $f[v] - f[u] = (1-w) - w = 1 - 2w$, the code's formula. $f[\text{root}]$ is by
construction the sum of per-edge non-compliances under the root-0 orientation --- the
number of reversals making $0$ a valid capital (each non-compliant edge needs exactly
one flip, and flipping precisely those yields the outward arborescence). By induction
along the preorder, every $f[v]$ equals the true reversal count for capital $v$; taking
the minimum over all $v$ and listing argmins gives exactly the problem's two output
lines (sorted by the label-order scan).
\end{proofbox}

\solstep{Implementation}
\begin{lstlisting}[caption={Codeforces 219D -- two-phase reroot on the reversal cost}, label={lst:c09-219d-soln}]
#include <bits/stdc++.h>
using namespace std;
int main() {
    int n;  scanf("%d", &n);
    vector<vector<pair<int,int>>> adj(n);      // (neighbour, disobey-cost one way)
    for (int i = 0; i < n - 1; i++) {
        int a, b;  scanf("%d %d", &a, &b);  --a; --b;
        adj[a].push_back({b, 0});              // road a->b: with it costs 0...
        adj[b].push_back({a, 1});              // ...against it costs 1
    }
    vector<int> par(n, -1), wIn(n), order;     // wIn[v] = cost of par[v]->v
    order.reserve(n);
    vector<int> st{0};  int f0 = 0;
    while (!st.empty()) {                       // phase 1: preorder (lst:c09-iter)
        int u = st.back();  st.pop_back();
        order.push_back(u);
        for (auto e : adj[u]) if (e.first != par[u]) {
            par[e.first] = u;  wIn[e.first] = e.second;
            f0 += e.second;                     // root 0 pays every wrong-way road
            st.push_back(e.first);
        }
    }
    vector<int> f(n);  f[0] = f0;  int best = f0;
    for (int u : order)                          // phase 2: shift along, forward
        for (auto e : adj[u]) if (e.first != par[u]) {
            f[e.first] = f[u] + 1 - 2 * wIn[e.first];   // +1 or -1 per flipped edge
            best = min(best, f[e.first]);
        }
    printf("%d\n", best);
    bool first = true;
    for (int v = 0; v < n; v++)                  // label order = ascending, free
        if (f[v] == best) { if (!first) printf(" "); first = false; printf("%d", v + 1); }
    printf("\n");
}
\end{lstlisting}

\solstep{Trace}
\begin{tracebox}[Official Sample 2: roads 1->4, 2->4, 3->4 (0-based: 0->3, 1->3, 2->3)]
Store adj[a].push(b,0), adj[b].push(a,1). Phase 1, preorder from node 0 (city 1):
push 3 via the COMPLIANT road 0->3 (cost 0); at 3, push 1 and 2 via stored 3->1, 3->2
cost 1 each (given roads read 1->3, 2->3). f[0] = 0+1+1 = \textbf{2} --- capital city~1
reverses the two roads it cannot reach along; its own road 1->4 is fine.\\
Phase 2 forward: f[3] = f[0] + 1 - 2*0 = 3 (capital 4: every road now points the wrong
way --- $3 > 2$, so 4 is not optimal, as the sample implies);\\
f[1] = f[3] + 1 - 2*1 = 2; \quad f[2] = f[3] + 1 - 2*1 = 2.\\
min = 2; argmins in label order = cities \textbf{1 2 3}: prints ``2 / 1 2 3''
\checkmark\ official. Sample 1 (roads 2->1, 2->3, i.e.\ node1->node0 and node1->node2;
start city 1 = node 0): 0's only neighbour is 1 through stored 0->1 cost 1 (against the
road); at 1, edge 1->2 is the given road, cost 0. f[0] = 1 + 0 = 1.\\
Phase 2: f[1] = f[0] + 1 - 2*1 = 0; \quad f[2] = f[1] + 1 - 2*0 = 1.\\
min 0 at city 2: prints ``0 / 2'' \checkmark.\end{tracebox}

\solstep{Complexity}
\begin{complexitybox}
Time: two linear phases, $\Theta(n)$ --- against $n \le 2\cdot 10^5$ this is the
required class, and the explicit stack keeps the chain case honest (CF stacks will not
hold $2\cdot 10^5$ frames by default). Memory: $\Theta(n)$. Verdict: \ac. The
$\pm 1$ shift you just hand-executed is the \emph{unweighted, binary-cost} instance of
the Summit's $+S - 2\,\mathit{mass}$ shift --- same skeleton, heavier payload.
\end{complexitybox}

\solstep{Pitfalls and Variants}
\begin{warningbox}
Failure modes with faces: reading the shift as $f[v] = f[u] \pm 1$ but picking the sign
from $v$'s \emph{new} edge cost instead of the \emph{down} cost $w$ (off by a
double: $+1-2w$ is $+1$ or $-1$ \emph{against the sign you first store}); sorting the
argmins by $f$-scan order instead of label order (Sample~2 prints \texttt{1 2 3}, not
\texttt{1 3 2} --- lucky here, a presentation-error elsewhere); and recursive DFS at
$2\cdot 10^5$. Variants: print the optimal capital's \emph{reversal list} (parent map
+ per-edge compliance); the weighted version where flipping costs money per road
(the $1-2w$ becomes a two-cost difference --- still one replay).
\end{warningbox}

\begin{insightbox}[Key takeaways]
\begin{enumerate}[nosep, leftmargin=1.6em]
  \item $\mathit{ecc}(v) = \max(d(v,A), d(v,B))$ from any diametral pair
        (\thmref{c09-eccend}) --- the radius/centre/whole-ecc-table cost two extra
        sweeps, not $n$ BFSs.
  \item The centre is the half-middle of every diameter, $R = \lceil D/2 \rceil$, one
        node or an adjacent pair (\thmref{c09-centerrule}) --- and never confused with
        the 12.3 centroid except in small examples.
  \item $\sum$-over-all-roots quantities move by a one-edge shift
        ($f[v] = f[u] + n - 2\size[v]$, or $+1-2w$ for arrow compliance): one DFS plus
        one replay reads every root --- the preview the Summit completes in weight.
\end{enumerate}
\end{insightbox}

\begin{hardeningbox}
\begin{itemize}[nosep, leftmargin=1.4em]
  \item \textbf{Brute-verify the identity.} For $v = 6$ and $v = 9$ of $T$, compute
        $\mathit{ecc}(v)$ by hand twice: all-node scan and $\max(d(v,8), d(v,10))$.
        Find the $u$ that attains each max and locate it on \figref{c09-eccprof}.
  \item \textbf{Shift both directions.} From $f(3) = 20$ in \exref{c09-ecc-run} derive
        $f(6)$ and $f(7)$ without re-summing (you will need $\size(6) = 1$).
  \item \textbf{Re-run 219D by eye} with root city~$4$ instead of~$1$: does the shift
        still land on the same argmin set? (It must --- root choice is bookkeeping.)
\end{itemize}
\end{hardeningbox}

\subsection{Practice Problems}

\begin{stuckbox}
\begin{itemize}[nosep, leftmargin=1.4em]
  \item Did your $\mathit{ecc}$ loop take the $\max$ over \emph{both} diameter-end
        distance arrays --- and do you still have the \emph{second} sweep's array
        around to index?
  \item On a two-node tree: one centre or two? ($D=1$ odd: \emph{two} adjacent
        centres --- run \thmref{c09-centerrule} and check your code agrees.)
  \item Did the 219D sign come out $+1$ for roads read with the flow and $-1$ against
        --- or the other way round? (Recheck with Sample~1's $f(\text{city }2)=0$.)
\end{itemize}
\end{stuckbox}

\begin{practiceset}
  \practice{Shortest Distance from All Buildings}{LeetCode 317}
    {https://leetcode.com/problems/shortest-distance-from-all-buildings/}{LC Hard · L5}
    {Grid with houses \texttt{1} and obstacles \texttt{2}: pick an empty land
    minimising the \emph{sum} of BFS distances to all houses; return $-1$ if none.
    The tree-free sibling of the $\sum$-family: BFS from \emph{each} house,
    accumulating \code{dist\[r\]\[c\] += d} and \code{reach\[r\]\[c\]++} on every cell
    the wavefront visits; answer $= \min\{\mathit{dist}[r][c] :
    \mathit{reach}[r][c] = \#\text{houses}\}$. $\bigO{H \cdot RC}$, which the
    constraints tolerate; the ``prune visited per house, never globally'' detail is
    where the WA lives. Premium-era page; statement verified live.}
\end{practiceset}

\medskip
\noindent\textbf{One to reread, not re-solve:} LC~310 \emph{Minimum Height Trees} is
already solved in Chapter~4 (\chref{c04}) by leaf-peeling -- go back and reread it with
this section's vocabulary: the one-or-two peeling survivors \emph{are} the centre
(\thmref{c09-centerrule}), and ``root the tree to minimise height'' is exactly ``root at
the centre.'' Prove to yourself the peeling and the diameter-middle algorithms always
agree --- no new code, just a new name for the object you already found.

\begin{recallbox}[10-second refresher -- 12.4 Eccentricity, Radius, Centre]
$\mathit{ecc}(v) = \max_u d(v,u) = \max(d(v,A), d(v,B))$ for any diameter ends $A,B$
(\thmref{c09-eccend}): three sweeps give the whole table. $R = \min_v
\mathit{ecc} = \lceil D/2 \rceil$; the centre is the middle node (even $D$) or adjacent
pair (odd $D$) of every diameter (\thmref{c09-centerrule}) --- $T$: $D=7$,
$\mathit{ecc}=[4,5,4,6,6,5,5,7,6,7]$, $R=4$, centre $\{1,3\}$. Centroid $\ne$ centre
(broom). Sums reroot by $f[v] = f[u] + n - 2\size[v]$ (or $\pm$ edge-cost
$\to$ \probref{c09-capital}, CF~219D) --- one DFS, one replay, every root.
\selftest{
  \item State \thmref{c09-eccend} and name the single claim about spurs it leans on.
  \item Why does $D = 7$ force \emph{two} centres, and $D = 6$ one? What is $R$ in each?
  \item Give the broom's centroid and centre; say which optimisation each answers.
  \item Derive the $\pm 1$ shift of \probref{c09-capital} from the $n - 2\size$ shift.
  \item A solver prints radius~$0$ for a long chain. Which trap did they fall into?
}
\end{recallbox}
